% Part: first-order-logic
% Chapter: beyond
% Section: modal-logics

\documentclass[../../../include/open-logic-section]{subfiles}

\begin{document}

\olfileid{fol}{byd}{mod}

\olsection{المنطقات الموجهة}

انظر في المثال الآتي لجملة شرطية:
\begin{quote}
  إذا كان جيريمي وحده في تلك الغرفة، فهو ثمل وعار ويرقص على الكراسي.
\end{quote}
هذا مثال لعبارة شرطية قد تكون صادقة ماديًا، لكنها مضللة مع ذلك، إذ
توحي بوجود صلة بين المقدمة والنتيجة أقوى من مجرد كون المقدمة كاذبة أو
التالية صادقة. أي إن الصياغة توحي بأن الدعوى ليست صادقة في هذا العالم
بعينه فحسب (وقد تكون صادقة فيه صدقًا تحصيليًا لأن جيريمي ليس وحده في
الغرفة)، بل بأن النتيجة \emph{كانت ستصدق} أيضًا \emph{لو} صدقت
المقدمة. وبعبارة أخرى، يمكن فهم العبارة على أن الدعوى تصدق لا في هذا
العالم وحده بل في كل عالم «ممكن»؛ أو أنها صادقة \emph{بالضرورة}، لا
صادقة في هذا العالم بعينه فحسب.

صُمم المنطق الموجه لتفسير هذا النوع من الضرورة. ونحصل على المنطق
القضوي الموجه من المنطق القضوي العادي بإضافة مؤثر الصندوق؛ أي إنه إذا
كانت $!A$~!!a{formula}، فإن $\Box !A$ كذلك. وتقرر $\Box !A$، حدسيًا،
أن $!A$ صادقة \emph{بالضرورة}، أو صادقة في كل عالم ممكن. ويؤخذ
$\Diamond !A$ عادة اختصارًا لـ~$\lnot \Box \lnot !A$، ويمكن قراءته
على أنه يقرر أن $!A$ صادقة \emph{بالإمكان}. ويمكن بالطبع إضافة الجهة
إلى منطق المحمولات أيضًا.

يمكن استعمال !!{structure}s كريبكه لتقديم دلالة للمنطق الموجه؛ بل إن
كريبكه صمم هذه الدلالة أولًا واضعًا المنطق الموجه نصب عينيه. وبدلًا من
الاقتصار على الترتيبات الجزئية، تكون لدينا بوجه أعم مجموعة من «العوالم
الممكنة»~$P$، وعلاقة «إمكان وصول» ثنائية $\Atom{R}{x,y}$ بين العوالم.
وتقرر $\Atom{R}{p,q}$، حدسيًا، أن العالم~$q$ متوافق مع~$p$؛ أي إنه إذا
كنا «في» العالم~$p$، فعلينا أن نقبل إمكان أن يكون العالم على صورة~$q$.

يسمى المنطق الموجه أحيانًا منطقًا «قصديًا»، في مقابل المنطق
«الامتدادي». فالدلالة المقصودة لمنطق امتدادي، مثل المنطق الكلاسيكي، لا
تحيل إلا إلى عالم واحد، هو العالم «الفعلي»؛ أما دلالة منطق «قصدي»
فتعتمد أنطولوجيا أغنى. وإلى جانب نمذجة الضرورة، يمكن استعمال الجهة
لنمذجة تراكيب لغوية أخرى، بإعادة تأويل $\Box$ و$\Diamond$ بحسب
التطبيق. فمثلًا:
\begin{enumerate}
\item في منطق القابلية للإثبات، يُقرأ $\Box !A$ على أنه «$!A$ قابلة
  للإثبات»، و$\Diamond !A$ على أنه «$!A$ متسقة».
\item في المنطق المعرفي يمكن قراءة $\Box !A$ على أنه «أعرف أن~$!A$»
  أو «أعتقد أن~$!A$».
\item في المنطق الزمني يمكن قراءة $\Box !A$ على أنه «$!A$ صادقة
  دائمًا»، و$\Diamond !A$ على أنه «$!A$ صادقة أحيانًا».
\end{enumerate}

نود أن نزيد المنطق قواعد وبديهيات تتناول الجهة. فالنسق~$\Log{S4}$،
مثلًا، يتألف من بديهيات المنطق القضوي وقواعده المعتادة، مع البديهيات
الآتية:
\begin{align*}
& \Box (!A \lif !B) \lif (\Box !A \lif \Box !B)\\
& \Box !A \lif !A \\
& \Box !A \lif \Box \Box !A
\intertext{ومع القاعدة «استنتج $\Box !A$ من~$!A$». ويضيف
  $\Log{S5}$ البديهية الآتية:}
& \Diamond !A \lif \Box \Diamond !A
\end{align*}
وقد تلائم تنويعات هذه البديهيات تطبيقات مختلفة؛ فيؤخذ S5 عادة، مثلًا،
على أنه يميز مفهوم الضرورة المنطقية. والميزة الحسنة هي أنه يمكن عادة
إيجاد دلالة يكون نسق ال!!{derivation} سليمًا وتامًا بالنسبة إليها،
بتقييد علاقة إمكان الوصول في !!{structure}s كريبكه بطرائق طبيعية.
فالنسق~$\Log{S4}$ يقابل، مثلًا، فئة !!{structure}s كريبكه التي تكون
فيها علاقة إمكان الوصول انعكاسية ومتعدية. ويقابل $\Log{S5}$ فئة
!!{structure}s كريبكه التي تكون فيها علاقة إمكان الوصول {\em كلية}،
أي التي يمكن فيها الوصول من كل عالم إلى كل عالم آخر؛ ومن ثم تصدق
$\Box !A$ إذا وفقط إذا صدقت $!A$ في كل عالم.

\end{document}
